\section{Model}
\label{sec:model}

The model is an individual-based, discrete-time stochastic process on a finite lattice. One time step
represents one week. Each individual carries four state variables: a position, an age, a sex and a
\emph{social state} $s_i\in[0,1]$, which summarises its capacity for normal social and reproductive
behaviour ($s=1$ fully competent, $s=0$ completely deteriorated). Individuals also keep a memory of
past contacts through a sparse affinity matrix. All rules are \emph{local}: no individual has access
to the population size or to any global statistic.

We first describe the \emph{core model} (Section~\ref{sec:model-baseline}): the lattice, the affinity matrix,
movement with attraction towards acquaintances, crowding damage to the social state, reproduction and senescent
mortality. Section~\ref{sec:model-rules} introduces four additional rules, R1, R2, AV and R7; each has a neutral
value that switches it off and leaves the core model unchanged. Section~\ref{sec:model-candidates} defines the four
calibrated candidates \textsf{C0}, \textsf{C1$'$}, \textsf{C1} and \textsf{C2} and the controls against which they are
compared, and Section~\ref{sec:model-groups} the dimensionless groups, the scales and the regimes we expect. The full
ODD description is in Appendix~\ref{app:odd}. Fig.~\ref{fig:model-schematic} summarises the ingredients and
Table~\ref{tab:params} lists all parameters.

\begin{figure*}[t]
  \centering
  \includegraphics[width=\textwidth]{model_schematic.pdf}
  \caption{Schematic of the model.
  \textbf{(a)}~A focal agent (ringed) chooses a cell of its Moore neighbourhood (dashed) with the weight of
  Eq.~\eqref{eq:model-move}; arrow width is proportional to the weight. Deteriorated individuals gather in pools
  ($m>m_c$) or occupy refuges (hatched, capacity $c_R=2$).
  \textbf{(b)}~Illustrative trajectories of $s$ versus age, from
  Eqs.~\eqref{eq:model-social}--\eqref{eq:model-learning} with fixed environments (not a simulation). A pup
  starts at $w\,s_{\rm mother}$ and gains competence only inside the window $a\le a_w$ (R1), at a rate set by
  its neighbours (R2).
  \textbf{(c)}~Phase diagram expected for \textsf{C0} (no refuges). Crowded and isolated classes coexist in
  the shaded band, $5.5\lesssim n^*=1/\kappa-1/\alpha\lesssim10$ with $\alpha\gtrsim3.5$. Points are
  calibration ensembles of \textsf{C0} labelled with $f_{\rm O17}$; filled markers have $f\ge0.8$. With
  refuges the band disappears (Q5, refuted; Section~\ref{sec:results}).}
  \label{fig:model-schematic}
\end{figure*}

\subsection{Core model}
\label{sec:model-baseline}

\paragraph{Space and affinity.}
Individuals live on an $L\times L$ lattice $G$ ($L=30$, $|G|=900$) with hard boundaries. Any number of
individuals may share a cell; $m_i$ denotes the number of occupants of the cell of $i$, itself included.
Every pair of individuals at Chebyshev distance $\le1$ at the end of a week reinforces its affinity,
$F_{ij}\leftarrow(1-\eta)F_{ij}+\eta$. All other pairs decay, $F_{ij}\leftarrow(1-\delta)F_{ij}$, and pairs
with $F_{ij}<\varepsilon$ are forgotten ($\eta=0.2$, $\delta=0.01$, $\varepsilon=10^{-6}$). The affinity
memory is therefore long, $1/\delta=100$ weeks, and is built only by proximity.

\paragraph{Movement.}
Every week each individual moves to one of the nine cells $y$ of its Moore neighbourhood, its own cell
included, cells outside $G$ excluded. All individuals move simultaneously, using the positions at the
start of the week. The probability of choosing $y$ is proportional to
\begin{equation}
  W_i(y)=\bigl[1+\alpha A_i(y)\bigr]\;
         \underbrace{e^{-\kappa(1-s_i)^{p}n_y}}_{\text{aversion}}\;
         \underbrace{\bigl[1+\pi(1-b_i)^{q}\,\mathbf 1_{\mathcal R}(y)\bigr]\;
         \bigl[1-\mathbf 1_{\mathcal F}(y)\bigr]}_{\text{refuges}},
  \qquad A_i(y)=\sum_{j\in y}F_{ij},
  \label{eq:model-move}
\end{equation}
where $n_y$ is the number of \emph{other} individuals in $y$. A refuge cell is \emph{full} if it already holds
$n_y\ge c_R$ other individuals; $\mathcal R$ is the set of non-full refuges and $\mathcal F$ the set of full ones, so
that $\mathbf 1_{\mathcal R}(y)=1$ if $y$ is a non-full refuge and 0 otherwise, and $\mathbf 1_{\mathcal F}(y)=1$ if
$y$ is a full refuge and 0 otherwise.
In the core model only the first factor is present ($\kappa=0$, no refuges, $\mathcal R=\mathcal F=\emptyset$): the parameter
$\alpha\ge0$ sets the attraction towards acquaintances. The aversion and refuge factors are the
additional rules of Section~\ref{sec:model-rules}.

\paragraph{Social state.}
After moving, the social state is updated with the local occupancy,
\begin{equation}
  s_i\leftarrow\min\Bigl\{1,\max\bigl\{0,\;s_i+\varrho_i-\gamma\,(m_i-m_c)_+\bigr\}\Bigr\},
  \label{eq:model-social}
\end{equation}
where $\gamma(m_i-m_c)_+$ is the damage inflicted by each occupant in excess of the community size
$m_c=8$ ($\gamma=0.1$). In the core model the recovery term is a constant, $\varrho_i=r$ ($r=0.01$).
Deterioration is thus a purely local, density-triggered process, and its convexity in the local
occupancy makes it act as a switch (Section~\ref{sec:meanfield}).

\paragraph{Reproduction.}
Every reproductive female ($a_{\min}\le a\le a_{\rm rep}$) that shares her cell with at least one
reproductive male chooses one of them uniformly at random, and conceives with probability
\begin{equation}
  p_{\rm conc}=p_0\,s_f\,s_m,\qquad p_0=0.3 .
  \label{eq:model-repro}
\end{equation}
The litter size $\ell\in\{3,\dots,10\}$ is drawn from a fixed distribution with mean $\bar\ell=6.55$.
Each pup survives birth with probability $q_0\,s_f$ ($q_0=0.95$), so deteriorated mothers lose their
young. Newborns are placed in the mother's cell, have a 1:1 sex ratio, and inherit
$s_{\rm nb}=w\,s_f+(1-w)\,s_0$. The core model has full inheritance, $w=1$; this is its central
hypothesis, that the lack of maternal care propagates to the offspring.
There is no gestation delay, and a female may give birth every week.

\paragraph{Mortality.}
Every individual dies at the end of the week with probability
\begin{equation}
  p_{\rm death}(a,s)=\Bigl[1+e^{-k\,(a-a_{\max}[1-\mu_s(1-s)])}\Bigr]^{-1},\qquad k=0.15 .
  \label{eq:model-mort}
\end{equation}
The factor $\mu_s\ge0$ would advance the senescence of deteriorated individuals. Every model in this paper,
candidates and controls, has $\mu_s=0$, so adult death is purely senescent. Calhoun attributes the decline to senescence and describes the beautiful ones as
``physically healthy''~\cite{Calhoun1973}. Anti-pattern A4 and the senescent character of the decline (O8) therefore
hold by construction: they are consistency checks of the design, not results.
There is no mortality before senescence. The population therefore cannot decrease before the first senescent deaths:
once natality stops it sits on a plateau, and $t_{\rm pk}=\arg\max N$ falls on the last birth that precedes the first
death (exactly so in 69\% of the runs of \textsf{C1}). A constant background hazard $\mu_{\rm bg}$, independent of age
and $s$, is a reported variant; it acts at every age, pups included. With $\mu_{\rm bg}=0.003$ per week the plateau
shrinks from 51 to 17 weeks and the primary outcome falls to 0.55 [0.48, 0.62], through O5 alone, which becomes
MARGINAL rather than failing (83 MARGINAL and 7 FAIL among the 90 runs that do not pass). Natality still ends at the
population maximum, $N(t_{B99})\approx0.99\,N_{\rm pk}$. What moves is the start of the plateau: deaths now balance the
last births, so $t'_{\rm pk}$ comes earlier and the last 1\% of the births falls more than $0.5\,t'_{\rm pk}$ after
it. The coexisting classes (O13, O13p), the patterns that discriminate against the null models, pass in every
run, and so does O11, also with $\mu_{\rm bg}=0.008$, where O5 fails in 70 of 200 runs
(Section~\ref{sec:res-robustness}).

\paragraph{Longevity, order of events, initial condition.}
The candidates use $a_{\max}=120$ and $a_{\rm rep}=80$ weeks (the short-lived control below uses 60 and 50). These
values match the menopause at about 560 days and the frequent longevity of about 800 days in Universe~25~\cite{Calhoun1973},
and they allow a decline governed by ageing (O8). Each week the order is: movement, affinities, update of
$s$, reproduction, ageing, mortality. The initial condition is dispersed: $N_0=200+200$
individuals of age 4 weeks with $s=1$, at random positions. Non-synchronous variants draw the initial ages
uniformly over 4--52 or 4--80 weeks; the primary outcome is then 0.93 and 0.63, again through O5 alone and by the
same mechanism (with 4--80 weeks, 73 of the 74 runs that do not pass are MARGINAL, with
$N(t_{B99})\ge0.97\,N_{\rm pk}$; Section~\ref{sec:res-robustness}). The founder protocol is in Appendix~\ref{sec:calib-founders}.

\subsection{Minimal additional rules}
\label{sec:model-rules}

The core model reproduces the explosive growth and a density-triggered loss of social state, and with the long
life span nearly every colony dies too (Section~\ref{sec:results}). With a short life span ($a_{\max}=60$) it
does not reproduce three facts of Universe~25: the irreversibility of the collapse, the failure of the young born
during the crowded phase, and the coexistence of distinct classes of deteriorated animals. With the long-lived
demography used as a null model the first two are largely there: the colony does not rebound (O6 and O7 in 0.99,
1.00 without the rules at $\alpha=4$) and, scored at the end of the plasticity window, its late cohorts fail too
(O11 in 0.81, and in 0.755 without the rules; Section~\ref{sec:res-null}). Only the coexistence of classes is missing. R1 is still needed
once refuges exist, because refuges create pockets of low density in which, without an irreversible window,
deteriorated animals recover and breed again: without R1 natality resumes (O6 and O7 in 0.22, O17p 0.21), and
\textsf{C1} without R1 and R2 also fails them (0.35, with O13p in 0.05), whereas without R1 and without refuges
the collapse stays irreversible in most runs (O6 and O7 in 0.85 of 40 runs without R1 and R7; 1.00 without any rule; Appendix~\ref{app:res-ablations}). R1 is
what lets a persistent isolated class coexist with an irreversible collapse. A fourth fact, the low peak with free space, needs no new rule. With the
weak attraction $\alpha=1.4$ of the long-lived control (Section~\ref{sec:model-candidates}) the peak reaches 0.85 of $K_\gamma$, with 59\% of the cells
empty, so O4 fails only through the peak density. With the calibrated $\alpha=4$ and none of the rules below, the peak
stays at 0.45 with 74\% of the cells empty, and the primary O4 passes in every run (Table~\ref{tab:res-null}). The
low peak with free space is therefore a property of strong attraction with crowding damage, not of the rules. We add one rule for each of the three facts (R1, R2, and AV or R7 for the classes),
justified by an observation of the original report~\cite{Calhoun1973}. Each rule has one
principal parameter, but not only one: R2 also fixes the partial inheritance $(w,s_0)$, and R7 has five parameters
($f_R$, $c_R$, $\pi$, $q$ and the basis $b$) plus the choice of a hard capacity, and was designed after \textsf{C0} had
failed to produce a persistent isolated class (Section~\ref{sec:circularity}). Whether a rule is \emph{necessary} depends on
the others. In \textsf{C0} each of R1, R2 and AV is necessary (Section~\ref{sec:calib-ablations}). With
refuges (R7), AV is no longer necessary for the primary outcome; it then only shapes the spatial
organisation (Section~\ref{sec:results} and the paragraph on AV below). The minimal set for the primary
outcome is therefore R1 + R2 + R7, which defines the candidate \textsf{C1$'$} (Section~\ref{sec:model-candidates}).
This minimality is conditional on the evaluation thresholds (class threshold 0.10, isolation $m\le2$, persistence
$\ge4$ weeks, ensemble threshold 0.8; Section~\ref{sec:res-ablations}): \textsf{C1$'$} passes the primary outcome in
0.835 [0.78, 0.88] of the runs, with a lower bound below the threshold.

\paragraph{R1: critical window ($a_w$).}
Recovery is only possible for $a\le a_w=12$ weeks: $\varrho_i=0$ for older individuals, while crowding
damage acts at every age. An adult is therefore a ratchet, $\dot s_i\le0$.
\emph{Motivation.} Behaviours that failed to develop in youth were not acquired later. Animals moved from
the mid-third of phase~D to new universes at very low density, ``even placing them with adequate sex
partners \dots\ that had matured in uncrowded conditions'', showed no recovery of reproductive behaviour
(H.~Marsden's work reported in~\cite{Calhoun1973}; see also~\cite{Marsden1972}). R1 makes a deteriorated
cohort unable to recover by individual means. Since conception scales as $s_fs_m$, an adult with $s\simeq0$ can
neither recover nor reproduce, whatever its partner. The failure of phase-D transplants (O12) is therefore a
consistency check of R1, which was motivated by the same observation, and not an independent result.

\paragraph{R2: juvenile socialisation by neighbours ($\lambda$, with partial inheritance $w$).}
Recovery is \emph{learned} from the social environment:
\begin{equation}
  \varrho_i=r\bigl[(1-\lambda)+\lambda\,\bar s_{V(i)}\bigr],\qquad
  \bar s_{V(i)}=\frac{1}{|V(i)|}\sum_{j\in V(i)}s_j ,
  \label{eq:model-learning}
\end{equation}
where $V(i)$ contains the other individuals of the $3\times3$ block centred on $i$, and
$\bar s_{V(i)}=0$ for an isolated individual. We use $\lambda=1$, $r=0.1$, and partial inheritance
$s_{\rm nb}=w\,s_f$ with $w=0.2$, $s_0=0$. With R1, Eq.~\eqref{eq:model-learning} acts only on juveniles:
a pup starts with a fraction of its mother's competence and must acquire the rest from competent
neighbours before age $a_w$.
\emph{Motivation.} ``By midway in phase~C essentially all young were prematurely rejected by their
mothers. They started independent life without having developed adequate affective
bonds''~\cite{Calhoun1973}. Competence is transmitted socially. Once competent adults become scarce,
pups have nobody to learn from, and the state $s=0$ is absorbing for the population, not just for an
individual. In \textsf{C1} this transmission is shallow: it fails in a single step. Crowding deteriorates the
founders during their own adult life, within about one generation time. Their offspring, about 70\% of all births,
are born below the threshold and learn little, because their neighbours are mostly other pups and founders that are
already deteriorating. The net reproductive number falls from 3.8 daughters per founder female to 0.4 in the first
generation born, and 97\% of the births belong to the first two generations (Appendix~\ref{app:calib-generations}). R2 does not
produce a cascade over many generations; it makes the failure of the first generations born during the boom
irreversible for the colony.
With $w=0.2$ and $s_0=0$ every pup is born with $s_{\rm nb}\le0.2$, that is, at or below the threshold that the
patterns use for ``deteriorated''. Part of the failure of the crowded-born cohorts is therefore definitional: with
$s_0=0.25$ pups are born at $s_{\rm nb}\approx0.4$, and O11, which reads the state at the age of maturity
($a_{\min}=6$ weeks), passes in only 8\% of the runs, so that the primary outcome falls to 0.08
(Section~\ref{sec:res-robustness}). The essential part of R2 is social
learning, not inheritance. With refuges, $w=0$ and $r=0.15$ give the primary outcome in every run
(Section~\ref{sec:results}). We keep $w=0.2$ for continuity with the inheritance hypothesis of the core model.

\paragraph{AV: aversion of the deteriorated ($\kappa$).}
Deteriorated individuals avoid occupied cells through the factor $e^{-\kappa(1-s_i)^{p}n_y}$ of
Eq.~\eqref{eq:model-move}, with $\kappa=0.15$. The exponent $p=4$ is fixed, so that the aversion only
matters for $s\lesssim0.3$. The attraction $\alpha$ does \emph{not} depend on $s$. For a deteriorated
individual facing a cell with $n$ acquaintances ($A\simeq n$), the product $(1+\alpha n)e^{-\kappa n}$ is
maximal at
\begin{equation}
  n^*=\frac{1}{\kappa}-\frac{1}{\alpha},
  \label{eq:model-nstar}
\end{equation}
which is $6.4$ for the calibrated values, while cells of strangers are always avoided. Deteriorated
individuals therefore split between groups of a few times $m_c$, where crowding damage continues, and
almost empty cells.
\emph{Motivation.} Calhoun distinguishes withdrawn males, ``aggregated in large pools near the centre of
the floor'' and wounded ``as a result of attacks by other withdrawn males'', from the beautiful ones,
which ``never engaged in fighting'', carried ``no wound or scar tissue'' and ``never attempted to
enter'' social life~\cite{Calhoun1973}. In the model the damage term $\gamma(m-m_c)_+$ inside a pool is
damage inflicted by other deteriorated individuals, and isolated individuals receive none (objective
O13). Making $\alpha$ itself decrease with $s$ instead homogenises the population and raises the peak
(Section~\ref{sec:calib}).
\emph{Status.} AV is necessary in \textsf{C0}, where it alone generates the isolated class. With refuges it
is not necessary for the primary outcome: without AV it passes in 0.835 [0.78, 0.88] of the runs, against
0.995 [0.97, 1.00] (Section~\ref{sec:calib-ablations}). AV then controls the use of space and the decline. It
\emph{reduces} the free space: without it the deteriorated animals gather in a few pools and half of the lattice is
empty at the peak ($f_0=0.52$), with it $f_0=0.16$ (Universe~25: about 20\% of nest boxes). It also makes the
decline senescent (O8 0.625 against 0.21). Neither enters the preregistered outcome, which is why the minimal
candidate \textsf{C1$'$} omits AV. The low peak with free space itself needs no rule (see above).

\paragraph{R7: refuges with preference and hard capacity ($\pi$).}
A fixed random fraction $f_R=0.3$ of the cells are refuges of capacity $c_R=2$. Nobody may enter, or
stay in, a refuge that already holds $c_R$ other individuals. The capacity is \emph{hard}: after each
movement step residents keep their place, entrants are admitted in random order up to the capacity, and the
rejected ones return to their previous cell; admission is iterated until no refuge exceeds its capacity, and pups
that would be born in a full refuge are placed in a neighbouring cell that is not a refuge. A simpler
single-pass admission, which does not count returning animals and newborns and lets a refuge hold up to 11
individuals, is a reported variant; it leaves the primary outcome unchanged (0.995 with either rule;
Section~\ref{sec:res-robustness}).
Deteriorated individuals prefer refuges through the factor $1+\pi(1-b_i)^{q}$ of
Eq.~\eqref{eq:model-move}, where $b_i$ is the state that drives the retreat. In \textsf{C1}, $b_i=s_i$ with
$q=4$ and $\pi=100$. The total refuge capacity is $f_Rc_R|G|=540$ places, about 20\% of the
adult population at the peak.
\emph{Motivation.} The Universe~25 enclosure offered 256 nest boxes, each holding about 15 mice, and
would not have limited shelter below 3840 animals. At the peak about 20\% of the nest boxes were
usually unoccupied, because many mice chose to crowd together. Withdrawn females retreated to the high,
less preferred nests~\cite{Calhoun1973}.
\emph{Mapping.} We use a single mapping. Every cell is a resting site, the analogue of a nest box, and $m_c$
plays the role of the comfortable occupancy (15 per box in Universe~25, $m_c=8$ in the model), so that
$K_\gamma=m_c|G|$ is the analogue of the 3840 places and $f_0$ of the fraction of empty boxes. The refuges are the
analogue of the high, less preferred boxes (the upper one or two of the four levels of each tunnel, 25--50\% of the
boxes, against $f_R=0.3$). Their capacity $c_R=2$ is \emph{not} that of a nest box: it is the isolation threshold
$m\le2$ of O13, a design choice (Section~\ref{sec:circularity}). Neither $|G|$ nor $m_c$ was calibrated to
Universe~25, so densities are compared with Calhoun's only in sign (Appendix~\ref{app:calib-scales}).

Both the preference and the hard capacity are necessary with synchronous moves. Without a preference ($\pi=0$) the
refuges stay empty, because a deteriorated individual surrounded by acquaintances prefers an occupied neighbouring
cell. With a soft capacity, simultaneous moves produce stampedes (Appendix~\ref{sec:calib-classes}). The hard
capacity is a remedy for the collisions of the synchronous update: with a random sequential update a soft capacity
also produces a persistent isolated class (O17p 0.995, $\rho_{\rm pk}=0.43$ against 0.36;
Section~\ref{sec:res-robustness}).

\paragraph{Variant of R7 in \textsf{C2}: retreat driven by maternal socialisation.}
In \textsf{C2} the retreat is driven by the competence of the mother at birth, $b_i=h_i\equiv s_{f}(t_{\rm birth})$,
with $h_i=1$ for founders, $q=1$ and $\pi=1000$. Pups of deteriorated mothers retreat, and
socialised adults that later deteriorate do not, so they end up crowded.
\emph{Motivation.} The beautiful ones belonged to ``the last half of the population born''; they were
young animals that never entered social life, not adults that withdrew~\cite{Calhoun1973}.
Basing the retreat on juvenile competence (the value of $s$ at age $a_w$) fails. Under R2 almost every
cohort born after the first ends the window with $s<0.1$, so there is no gradient between cohorts.

\paragraph{Rules tested and rejected.}
Several rules were unnecessary or destroyed some pattern, and are not part of any candidate:
\begin{itemize}
  \item long-range jumps;
  \item an attraction that decreases with $s$;
  \item territoriality and defeat of males;
  \item social mortality;
  \item bonds restricted to competent pairs;
  \item movement sub-steps.
\end{itemize}
Mating in the Moore neighbourhood (\texttt{mate\_radius}${}=1$) only matters for the founder protocol.

\subsection{Calibrated candidates}
\label{sec:model-candidates}

We consider four candidates:
\begin{description}
  \item[\textsf{C0}] core model + longevity + R1 + R2 + AV. This is the minimal model that satisfies the essential
    patterns evaluable with $N_0=400$ when O13 is read instantaneously (O17 in 0.92 of the runs).
  \item[\textsf{C1$'$}] core model + longevity + R1 + R2 + R7 ($b=s$), without AV. This is the minimal set
    of rules for the preregistered primary outcome O17p (all essential patterns and a \emph{persistent}
    isolated class). See Section~\ref{sec:calib-ablations}.
  \item[\textsf{C1}] \textsf{C1$'$} + AV, equivalently \textsf{C0} + R7. AV adds the spatial organisation: it spreads
    the deteriorated animals over the lattice (the empty fraction at the peak falls from 0.52 to 0.16) and makes the
    decline senescent. \textsf{C1} is the reference of the production design.
  \item[\textsf{C2}] \textsf{C0} + R7 with the retreat driven by maternal socialisation ($b=h$). It makes the isolated
    class a later-born cohort.
\end{description}
The presets are \texttt{calhoun\_C0}, \texttt{calhoun\_C1p}, \texttt{calhoun\_C1} and \texttt{calhoun\_C2},
and Table~\ref{tab:params} lists their parameters.

\paragraph{Controls.}
A rule is switched off by its neutral value: $a_w=\infty$ for R1; $\lambda=0$ with full inheritance ($w=1$) and
individual recovery at the core rate $r=0.01$ for R2; $\kappa=0$ for AV; no refuge cells ($f_R=0$) for R7. Two
controls are the core model alone, with all four rules switched off, attraction $\alpha=1.4$ and the dispersed
initial condition:
\begin{description}
  \item[Short-lived control] a short life span, $a_{\max}=60$ and $a_{\rm rep}=50$ weeks;
  \item[Long-lived control] the life span of the candidates, $a_{\max}=120$ and $a_{\rm rep}=80$ weeks.
\end{description}
Table~\ref{tab:params} gives both in full. The null models of Section~\ref{sec:res-null} are run in a separate
ensemble and are defined relative to \textsf{C1}: no crowding damage ($\gamma=0$); $-$R2$^*$, \textsf{C1} with R2
switched off; the long-lived control; and ``no rules'', \textsf{C1} with all four rules switched off, which is the
long-lived control with the calibrated attraction $\alpha=4$. Two partial ablations of \textsf{C1} complete the
comparison: $\alpha=0$ (no attraction) and \textsf{C0} (no refuges).

\begin{table}[t]
  \centering
  \caption{Parameters of the short-lived (SL) and long-lived (LL) controls, which are the core model with the
  four rules switched off, and of the calibrated candidates \textsf{C0}, \textsf{C1$'$}, \textsf{C1} and \textsf{C2}.
  Time in weeks. A dash means that the rule is inactive. The code name is the field of \texttt{Params} in
  the \texttt{u25} package.}
  \label{tab:params}
  \small
  \resizebox{\textwidth}{!}{%
  \begin{tabular}{@{}lllcccccc@{}}
    \toprule
    Symbol & Meaning & Code name & SL & LL & \textsf{C0} & \textsf{C1$'$} & \textsf{C1} & \textsf{C2} \\
    \midrule
    \multicolumn{9}{@{}l}{\emph{Space, motion and affinity}}\\
    $L$ & lattice side (hard walls) & \texttt{grid\_size} & 30 & 30 & 30 & 30 & 30 & 30 \\
    $\alpha$ & attraction to acquaintances & \texttt{alpha} & 1.4 & 1.4 & 4.0 & 4.0 & 4.0 & 4.0 \\
    $\eta,\ \delta$ & affinity reinforcement, decay & \texttt{eta}, \texttt{delta} & 0.2, 0.01 & 0.2, 0.01 & 0.2, 0.01 & 0.2, 0.01 & 0.2, 0.01 & 0.2, 0.01 \\
    $\kappa,\ p$ & aversion of the deteriorated & \texttt{crowd\_aversion}, \texttt{\_power} & -- & -- & 0.15, 4 & -- & 0.15, 4 & 0.15, 4 \\
    \midrule
    \multicolumn{9}{@{}l}{\emph{Social state}}\\
    $m_c$ & community size & \texttt{m\_c} & 8 & 8 & 8 & 8 & 8 & 8 \\
    $\gamma$ & damage per excess occupant & \texttt{gamma} & 0.1 & 0.1 & 0.1 & 0.1 & 0.1 & 0.1 \\
    $r$ & recovery / learning rate & \texttt{rec} & 0.01 & 0.01 & 0.1 & 0.1 & 0.1 & 0.1 \\
    $a_w$ & R1 critical window & \texttt{plast\_age\_max} & $\infty$ & $\infty$ & 12 & 12 & 12 & 12 \\
    $\lambda$ & R2 weight of neighbours & \texttt{social\_rec\_lambda} & 0 & 0 & 1 & 1 & 1 & 1 \\
    $w,\ s_0$ & inheritance $s_{\rm nb}=w s_f+(1-w)s_0$ & \texttt{inherit\_weight}, \texttt{s\_newborn\_base} & 1, -- & 1, -- & 0.2, 0 & 0.2, 0 & 0.2, 0 & 0.2, 0 \\
    \midrule
    \multicolumn{9}{@{}l}{\emph{Refuges}}\\
    $f_R,\ c_R$ & refuge fraction, capacity (hard) & \texttt{refuge\_frac}, \texttt{\_capacity} & -- & -- & -- & 0.3, 2 & 0.3, 2 & 0.3, 2 \\
    $\pi,\ q$ & refuge preference & \texttt{refuge\_pref}, \texttt{\_power} & -- & -- & -- & 100, 4 & 100, 4 & 1000, 1 \\
    $b_i$ & state driving the retreat & \texttt{refuge\_pref\_basis} & -- & -- & -- & $s_i$ & $s_i$ & $h_i$ (mother) \\
    \midrule
    \multicolumn{9}{@{}l}{\emph{Demography}}\\
    $N_0$ & initial individuals (age 4, $s=1$) & \texttt{n\_males}+\texttt{n\_females} & 400 & 400 & 400 & 400 & 400 & 400 \\
    $p_0$ & conception probability & \texttt{p0} & 0.3 & 0.3 & 0.3 & 0.3 & 0.3 & 0.3 \\
    $q_0,\ \bar\ell$ & neonatal survival, mean litter & \texttt{neonatal\_survival} & 0.95, 6.55 & 0.95, 6.55 & 0.95, 6.55 & 0.95, 6.55 & 0.95, 6.55 & 0.95, 6.55 \\
    $a_{\min},\ a_{\rm rep}$ & reproductive ages & \texttt{repro\_age\_min}, \texttt{\_max} & 6, 50 & 6, 80 & 6, 80 & 6, 80 & 6, 80 & 6, 80 \\
    $a_{\max},\ k$ & senescence age, steepness & \texttt{a\_max}, \texttt{k\_mort} & 60, 0.15 & 120, 0.15 & 120, 0.15 & 120, 0.15 & 120, 0.15 & 120, 0.15 \\
    $\mu_s$ & social advance of senescence & \texttt{social\_mortality\_factor} & 0 & 0 & 0 & 0 & 0 & 0 \\
    \bottomrule
  \end{tabular}}
\end{table}

\subsection{Dimensionless groups, scales and expected regimes}
\label{sec:model-groups}

The qualitative behaviour depends on a few combinations of parameters; the values quoted are for \textsf{C1}. The
scales of the model and their counterparts in Universe~25 are listed in Appendix~\ref{app:calib-scales}.

\begin{itemize}
  \item $\Pi_\gamma=\gamma/r=1$: crowding damage per excess neighbour relative to the weekly learning rate.
  \item $\Lambda=r\,a_w=1.2$: the largest competence a juvenile can learn inside the window. A socialised
    lineage sustains itself only if $w+\Lambda\langle\bar s_V\rangle\gtrsim1$. Fast growth lowers the
    competence of the neighbours $\langle\bar s_V\rangle$, because novices outnumber teachers (a
    \emph{learning dilution}). Without crowding the colony goes extinct below a boundary in $(w,\Lambda)$ and grows
    without bound at the calibrated values (Q7). When $a_w$ is varied as well, $r$ and $a_w$ do not act through
    their product alone: the extinction transition is located only to within the spacing of the grid, and the local
    weight of $a_w$ relative to $r$ is about one half between $a_w=6$ and 12 weeks and about one between 12 and 24
    (Section~\ref{sec:res-Q}). Crowding is the trigger and R2 the amplifier.
  \item $n^*/m_c\approx0.8$ (Eq.~\eqref{eq:model-nstar}): the preferred group size of a deteriorated individual
    relative to the damage threshold. Without refuges it decides between pools, isolation, or both; with
    refuges it only marks the onset of homogenisation, near $n^*\approx3$--$4$ (Q5).
  \item $\phi_R=f_Rc_R|G|/N_{\rm ad}(t_{\rm pk})\approx0.2$: refuge capacity per adult at the peak. It bounds the
    persistently isolated fraction, $\Phi_{\rm BO}^{\rm pers}\le\phi_R$. The bound holds at all 50 points of the
    refuge plane, so a class of $\ge10\%$ needs $f_Rc_R\gtrsim0.3$; in that plane a persistent class appears only for
    $f_Rc_R\ge0.4$ and a strong preference (Section~\ref{sec:res-robustness}).
  \item $n_0=N_0/|G|=0.44$ and the Damk\"ohler number $\mathrm{Da}=t_{\rm mix}/t_{\rm det}$. Here
    $t_{\rm mix}=L^2/(4D_{\rm eff})$ is the time an individual needs to move across the lattice, with $D_{\rm eff}$ the
    mean squared displacement over 16 weeks divided by $4\times16$, and $t_{\rm det}=t_{1/2}-t_{\rm on}$ is the time
    the animals past the critical window (age $\ge a_w$) need to lose half of their competence ($t_{1/2}$, the first
    week with their mean $s<0.5$) after crowding sets in ($t_{\rm on}$, the first week with 5\% of the individuals in
    cells with $m>m_c$). These animals no longer learn, and newborns do not enter their mean until they leave the
    window, so $t_{\rm det}$ measures deterioration and not the dilution by pups born with $s\le0.2$. In \textsf{C1}
    $D_{\rm eff}\approx0.13$ cells$^2$ per week (8 exploratory runs), $t_{\rm mix}\approx1.7\times10^3$ weeks and
    $t_{\rm det}\approx10$ weeks ($t_{\rm on}=6$, $t_{1/2}=16$; reference ensemble), so
    $\mathrm{Da}\approx1.7\times10^2$. Deterioration is correlated over
    $\xi=(4D_{\rm eff}t_{\rm det})^{1/2}\approx2$ cells, the size of the contact neighbourhood. In founder colonies the
    relevant transport is the colonisation front: the filling time $t_{\rm fill}\approx t_{\rm pk}\approx170$--215
    weeks against $t_{\rm det}\approx20$--30 gives $\mathrm{Da}_{\rm fill}\approx10$
    (Appendix~\ref{sec:calib-founders}). In Universe~25 mice crossed the enclosure daily and the social organisation
    took a whole phase~C, about 35 weeks, to die, so $\mathrm{Da}\lesssim4\times10^{-3}$. The spatial structure of the
    model therefore emerges among animals that are almost sessile on the time scale of deterioration, which is not
    Calhoun's regime. Colonies grow from local nuclei, so $\rho_{\rm pk}$ falls with $n_0$, from 0.40 at $N_0=400$ to
    0.03 for eight founders.
  \item $\Gamma=t_{\rm det}/T_{\rm gen}\approx1.1$: deterioration time relative to the generation time,
    $T_{\rm gen}\approx9$ weeks (the interval between the median birth weeks of the first and second generations);
    the interquartile range over runs is 1.0--1.29, and the founders, a closed cohort, give 1.2. The mean over all
    animals of age $\ge a_{\min}=6$ would give 0.8, but that mean is diluted by the first generation, which
    enters it with $s\le0.2$ while still inside the window: when it falls below 0.5, in week 13, the founders still
    have $\bar s=0.66$ and are 39\% of those adults. Deterioration therefore takes of the order of one generation.
    The value depends on how $T_{\rm gen}$ is defined: with overlapping cohorts the interval between median birth
    weeks is not a demographic generation time, and in the null models it drops to 5 weeks (4 in some runs), below
    the age of first reproduction, which inflates their $\Gamma$ (1.4--1.6). With the standard demographic
    definition, the mean age of the mothers at birth ($T_{\rm gen}=14.3$ weeks in \textsf{C1}), $\Gamma\approx0.7$
    in \textsf{C1} and 0.60--0.85 in the other arms where it was computed; $\Gamma$ is of order one under either
    definition (Appendix~\ref{app:calib-generations}).
    $\Gamma$ does not set the depth of the boom. With five times less crowding damage ($\gamma=0.02$) $\Gamma=1.6$
    (1.8 on the founders), and the boom is still one to two generations deep: the mean generation of the births is
    1.42 against 1.35, $R_{\rm net}$ of the first generation is 0.52, and O17p passes in 17 of 20 runs. Pups born
    at $s_{\rm nb}\approx0.4$ ($s_0=0.25$) deepen the boom more at about the same $\Gamma$ (1.3): the mean
    generation is 1.51 and $R_{\rm net}=0.62$ (20 runs per point;
    Appendix~\ref{app:calib-generations}). The depth is set in the first step, from the founders to their offspring:
    pups are born below the deterioration threshold, conception and pup survival scale with the parents' state, and
    pups learn from neighbourhoods made mostly of other pups and of founders that are already deteriorating. In
    Universe~25 the deterioration took the whole phase~C, about 35 weeks, with $T_{\rm gen}\approx10$ weeks, so
    $\Gamma\approx3.5$ (our estimate). With a different operational definition, that is an order-of-magnitude
    comparison, not a test.
  \item \emph{Extinction time.} Without mortality before senescence, $T_{\rm ext}=\max_i(b_i+\ell_i)\simeq
    t_{\rm last}+\ell_{\max}$: the time of the last birth plus the longest residual lifetime. In \textsf{C1}
    $\ell_{\max}\approx110$ weeks with an interquartile range of a few weeks, of the order of the inverse slope of the
    logistic hazard, $1/k\approx7$ weeks. The spread of $T_{\rm ext}$ is therefore that of $t_{\rm last}$, the
    extreme of a sparse tail of births whose rate decays over $\tau_B=(t_{B99}-t_{B90})/\ln10\approx7$ weeks
    ($t_{Bq}$ is the week by which a fraction $q$ of all births has occurred). At fixed density the median of
    $T_{\rm ext}$ grows with the logarithm of the system size, from 162 weeks at $L=15$ to 184 at $L=60$, and its
    interquartile range over the median is 0.15 for $L\le20$ and about 0.08 for $L\ge30$
    (Appendix~\ref{app:calib-evs}). A small dispersion of $T_{\rm ext}$ is a property of how abruptly natality stops
    and of the size of the system, not a signature of the rules. In the same way
    $(T_{\rm ext}-t_{\rm pk})/t_{\rm pk}\approx2.1$ (Universe~25: $\ge1.84$) depends on where $t_{\rm pk}$ falls on the
    plateau: it is 4.0 if the peak time is the start of the plateau, $t'_{\rm pk}$, and the plateau itself shrinks
    from 51 to 17 weeks with a background mortality of 0.003 per week. Similarly $t_{\rm pk}/a_{\max}\approx0.45$
    (0.28 with $t'_{\rm pk}$), against 0.7 in Universe~25.
\end{itemize}

\paragraph{Expected regimes and their test.}
Before the production sweeps, the mean field (Section~\ref{sec:meanfield}) and the calibration suggested four
regimes. The production planes test them (Section~\ref{sec:results}). The planes report the fraction of runs that
pass a composite classifier (O17p). That fraction is not an order parameter, and its boundaries mix the mechanism
with the thresholds of the definitions. We therefore also report continuous order parameters: $\rho_{\rm pk}$, the
empty fraction $f_0$, the class order parameter $\Psi=\min(\Phi_{\rm BO}^{\rm pers},\Phi_{\rm CR})$ (O13p is
$\Psi\ge0.1$), the net reproductive number of the first generation and $\tau_B$ (Section~\ref{sec:res-robustness}; the widths of the
transitions could not be resolved with 20--30 runs per point).
\begin{enumerate}
  \item \emph{Recurrent growth} (A2), when recovery is individual and reversible. \emph{Confirmed}: without
    the window $P_{\rm ext}$ falls from 1 at $r=0.02$ to 0 for $r\ge0.2$.
  \item \emph{Unbounded growth} (A1), when crowding damage is weak or the population is homogenised.
    \emph{Confirmed} for $\gamma=0$.
  \item \emph{A single collapse without recovery}, the Calhoun-like regime. Extinction alone does not
    identify it, since the long-lived demography without new rules also goes extinct. What identifies it is
    the structure of the collapse, which requires R1 and R2. In \textsf{C0} the classes are set by
    $(\alpha,\kappa)$, inside the band of Fig.~\ref{fig:model-schematic}(c). With refuges the band
    disappears (Q5, refuted). No recovery appears in larger systems either: at fixed density, A2 is cleared in all
    150 exploratory runs with $L$ from 20 to 84 (30 per size, with the single-pass refuge admission), up to about
    21\,000 animals (Appendix~\ref{app:calib-evs}).
  \item \emph{Trivial early collapse}: founders in one cell, where local nuclei deteriorate before the colony
    spreads ($\mathrm{Da}_{\rm fill}\approx10$), or very strong aggregation.
\end{enumerate}
Persistence of the isolated class is a separate axis: it requires $\phi_R\gtrsim0.1$ and $\pi\gtrsim100$, and
its cohort structure depends on the basis $b$ of the retreat.
